\title{Scaling Language Models: Methods, Analysis & Insights from Training Gopher}

\begin{abstract}
Language modelling provides a step towards intelligent communication systems by harnessing large repositories of written human knowledge to better predict and understand the world. In this paper, we present an analysis of Transformer-based language model performance across a wide range of model scales --- from models with tens of millions of parameters up to a 280 billion parameter model called \textit{Gopher}. These models are evaluated on 152 diverse tasks, achieving state-of-the-art performance across the majority. Gains from scale are largest in areas such as reading comprehension, fact-checking, and the identification of toxic language, but logical and mathematical reasoning see less benefit. We provide a holistic analysis of the training dataset and model’s behaviour, covering the intersection of model scale with bias and toxicity.  Finally we discuss the application of language models to AI safety and the mitigation of downstream harms.
\end{abstract}

\section{Introduction}
Natural language communication is core to intelligence, as it allows ideas to be efficiently shared between humans or artificially intelligent systems. The generality of language allows us to express many intelligence tasks as taking in natural language input and producing natural language output.

Autoregressive language modelling --- predicting the future of a text sequence from its past --- provides a simple yet powerful objective that admits formulation of numerous cognitive tasks. At the same time, it opens the door to plentiful training data: the internet, books, articles, code, and other writing. However this training objective is only an approximation to any specific goal or application, since we predict everything in the sequence rather than only the aspects we care about. Yet if we treat the resulting models with appropriate caution, we believe they will be a powerful tool to capture some of the richness of human intelligence.

Using language models as an ingredient towards intelligence contrasts with their original application: transferring text over a limited-bandwidth communication channel. Shannon's Mathematical Theory of Communication~\citep{shannon1948mathematical} linked the statistical modelling of natural language with compression, showing that measuring the cross entropy of a language model is equivalent to measuring its compression rate. Shannon fit early language models to real data via precomputed tables of text statistics~\citep{dewey1923relativ} relating model complexity to improved text compression alongside more realistic text generation.\footnote{A sample from Shannon's word-pair model: ``the head and in frontal attack on an english writer that the character of this point is therefore another method for the letters that the time of who ever told the problem for an unexpected.''}  But the relation to intelligence was there from the start: Shannon posits that a sufficiently complex model will resemble human communication adequately, and the \textit{Imitation Game}~\citep{turing1950computing} cemented the link. The relation between data compression (via prediction) and intelligence has been further expanded upon since (see \citet{wolff1982language, chater1999search, legg2007universal}).

A key driver towards better language models has been modern computing. From their pen-and-paper origins, language models have transformed in capacity and predictive power by the exponential rise in compute~\citep{moore1965cramming}. In the 1990s and 2000s, $n$-gram models saw increases in scale and better smoothing approaches~\citep{ney1994structuring}, including a 300 billion $n$-gram model trained on two trillion tokens of text~\citep{brants2007large}.  These models have been applied to speech recognition~\citep{jelinek1997statistical}, spelling correction~\citep{brill2000improved}, machine translation~\citep{brown1990statistical}, and many other areas. However $n$-gram models become statistically and computationally inefficient as the context length is increased, which limits the richness of language they can model.

In the past two decades language models have progressed to neural networks that capture the structure of language implicitly~\citep{bengio2003neural, mikolov2010recurrent, graves2013generating, jozefowicz2016exploring, radford2019language}. Progress has been driven by both scale and network architecture~\citep{hochreiter1997long,bahdanau2014neural,vaswani2017attention}.
\citet{rosenfeld2020a} and \citet{kaplan2020scaling} independently found power laws relating cross entropy loss to model size for recurrent and Transformer neural language models respectively.  The empirically predicted gains to scale were realised in practice by the \textit{Generative Pre-trained Transformer 3} (GPT-3, \citet{gpt3}), a 175 billion parameter Transformer trained over 300 billion tokens of text, which consumed zettaflops of compute to train --- an order of magnitude beyond prior work~\citep{rosset2020turing}. GPT-3 demonstrated unprecedented generation quality alongside generalist capabilities across many Natural Language Processing (NLP) tasks --- notably when prompted with examples (termed few-shot prompting).

In this paper we describe a protocol for training a state-of-the-art large language model and present a 280 billion parameter model called \textit{Gopher}. We outline the methods of architecture specification, optimisation, infrastructure, and the curation of a high-quality text dataset \textit{MassiveText} in \autoref{sec:method}. We perform a broad analysis of benchmark performance across 152 tasks that examine several diverse aspects of intelligence, and summarise the key results in \autoref{sec:results}. We see that \textit{Gopher} lifts the performance over current state-of-the-art language models across roughly 81\% of tasks containing comparable results, notably in knowledge-intensive domains such as fact checking and general knowledge. 

As harmful content occurs both in \textit{Gopher}'s training set and in many potential downstream applications, we examine model toxicity and bias in~\autoref{sec:model_analysis} with a focus on how scale influences these properties. We find larger models are more likely to generate toxic responses when provided with toxic prompts, but they can also more accurately classify toxicity. We also analyse \textit{Gopher} in a dialogue-interaction setting in \autoref{sec:dialogue} via prompting and present several transcripts to demonstrate qualitative capabilities and limitations of the model.

Finally, we discuss the ethical and safe application of these models including which types of undesirable behaviour to mitigate before and after training in \autoref{sec:discussion}. We discuss application-driven safety and the potential for language models to accelerate research towards safer intelligent technology.

\section{Background}
\label{background}
Language modelling refers to modelling the probability of text $P(S)$ where $S$ can be a sentence, paragraph, or document depending on the application. This is done by \textit{tokenizing} the string: mapping it to a sequence of integer-valued \textit{tokens}: $g(S) = X = (X_1, X_2, \ldots, X_n) \in V^n$ where $V$ is the vocabulary (a finite set of positive integers) and $n$ is the resulting sequence length, and modelling $X$. Tokenization can be \textit{open-vocabulary} where any string can be uniquely tokenized, e.g., byte-level modelling, or \textit{closed-vocabulary} where only a subset of text can be uniquely represented, e.g., a list of words and a singular out-of-vocabulary token. We employ open-vocabulary tokenization via a mixture of byte-pair encoding (BPE) with a backoff to UTF-8 bytes in the style of~\citet{radford2018improving}. 

The typical way to model the token sequence $X$ is via the \textit{chain rule} $P(X) = P(X_1, X_2, \ldots, X_n) = \prod_{i=1}^n P(X_i | X_{<i})$. This is also known as \textit{autoregressive} sequence modelling, because at each time-step the future (in this case, future token) is predicted based upon the past context. Whilst there are other objectives towards modelling a sequence, such as modelling masked tokens given bi-directional context~\citep{mikolov2013efficient, devlin2019bert} and modelling all permutations of the sequence~\citep{yang2019xlnet} we focus on autoregressive modelling due to its strong performance and simplicity. We shall refer to language models hereon as the function approximators to perform next-token prediction.

A class of neural networks known as Transformers~\citep{vaswani2017attention} have demonstrated state-of-the-art language model performance in recent years~\citep{dai2019transformer, radford2018improving, radford2019language} and this is the architecture we focus on in this paper. There has been a trend of scaling the combination of training data, model size (measured in parameters) and training computation to obtain models with improved performance across academic and industrial benchmarks. Notable models along this progression include the 345 million parameter BERT~\citep{devlin2019bert} performing strongly across a wide benchmark of language classification tasks, the 1.5 billion parameter GPT-2~\citep{radford2018improving} and 8.3 billion parameter Megatron ~\citep{shoeybi2019megatron} displaying progressively superior zero-shot language model performance, the 11 billion parameter T5~\citep{raffel2019exploring} which advanced transfer learning and performance on several closed-book question answering tasks, and the aforementioned 175 billion parameter GPT-3. The moniker \textit{Large Language Models} (LLMs) has become popular to describe this generation of larger models.

Since GPT-3 there has been a 178B parameter Transformer language model Jurassic-1~\citep{jurassic} which uses a diverse training set and a larger tokenizer vocabulary size, along with an announced 530B Megatron-Turing NLG~\citep{Megatron-Turing} which trains on a released dataset (The Pile, ~\citet{pile}) (which we evaluate on) and has reported some tentative performance numbers. There have also been Transformer variants which incorporate a sparse mixture of experts ~\citep{fedus2021switch, roller2021hash} to increase the model size (in some cases to trillions of parameters) with more modest compute budgets. Other recent LLMs include two models (FLAN and T0) fine-tuned on instructions for an array of down-stream tasks ~\citep{sanh2021multitask, wei2021finetuned} which improves performance to unseen tasks --- these ideas are complementary to the initial task of building a powerful language model but we compare performance nonetheless where possible.

\section{Method}
\label{sec:method}

\subsection{Models}
\label{method:models}
In this paper we present results on six Transformer language models ranging from 44 million to 280 billion parameters, with the architectural details displayed in \autoref{tab:arch}. We refer to the largest as \textit{Gopher} and the entire set of models as the \textit{Gopher} family.

We use the autoregressive Transformer architecture detailed in \citet{radford2019language} with two modifications: we use RMSNorm \citep{zhang2019root} instead of LayerNorm~\citep{ba2016layer}, and we use the relative positional encoding scheme from \citet{dai2019transformer} rather than absolute positional encodings. Relative encodings permit us to evaluate on longer sequences than we trained on, which improves the modelling of articles and books as shown in \autoref{appendix:context_len_scaling}.
We tokenize the text using SentencePiece~\citep{kudo2018sentencepiece} with a vocabulary of 32,000 and use a byte-level backoff to support open-vocabulary modelling. The \textit{Gopher} model card \citep{mitchell2019model} is included in \autoref{appendix:gopher-model-card}.

\subsection{Training}
\label{method:training}
We train all models for 300 billion tokens with a 2048 token context window, using the Adam \citep{kingma2014adam} optimiser. We warm-up the learning rate from $10^{-7}$ to the maximum learning rate over the first 1500 steps, and then decay it 10$\times$ using a cosine schedule. As we increase model size, we decrease the maximum learning rate and increase the number of tokens in each batch, as shown in \autoref{tab:arch}.
Furthermore, we increase \textit{Gopher}'s batch size from three to six million tokens per batch during training. We clip gradients based on the global gradient norm using a clipping value of 1. However, for the 7.1B model and for \textit{Gopher} we reduce this to 0.25 for improved stability.

We incorporate the \texttt{bfloat16} numerical format to reduce memory and increase training throughput. Models smaller than 7.1B are trained with mixed precision \texttt{float32} parameters and \texttt{bfloat16} activations \citep{micikevicius2017mixed}, while 7.1B and 280B use \texttt{bfloat16} activations \textit{and} parameters. \texttt{bfloat16} parameters are updated using stochastic rounding to maintain stability~\citep{gupta2015deep}. We subsequently found that stochastic rounding does not fully recover mixed precision training performance; more details can be found in \autoref{app:lessons_learned}.

\subsection{Infrastructure}
\label{method:infra}
We built our training and evaluation codebase with JAX \citep{jax2018github} and Haiku \citep{haiku2020github}. In particular, we use JAX's \texttt{pmap} transformation to efficiently express both data and model parallelism. We trained and evaluated all models on TPUv3 chips~\citep{tpuacm}.

The half-precision parameters and single-precision Adam state for \textit{Gopher} occupy 2.5 TiB, which far exceeds the 16 GiB of memory available on each TPUv3 core. To address these memory concerns, we use optimiser state partitioning \citep{rajbhandari2020zero}, model parallelism \citep{shoeybi2019megatron}, and rematerialisation \citep{griewank2000algorithm} to partition the model state and reduce the activations so that they fit in TPU memory.

We find that both data and model parallelism are low-overhead on TPUv3s due to their fast cross-chip communication and only incur a 10\% overhead when training \textit{Gopher}. Therefore, we find that pipelining \citep{huang2019gpipe} is not necessary on TPUs until the training scale exceeds the 1024-chip ``pod'', which greatly simplifies training mid-sized models. However, pipelining is an efficient parallelism method on commodity networks due to its low communication volume, so is well suited to connecting multiple TPU pods. In summary, we train \textit{Gopher} by using model and data parallelism within TPU pods and pipelining across them. We verified through simulation that this topology was sensible for our hardware \citep{automap2021}; see \autoref{tab:scaling_time} for details.

\subsection{Training Dataset}
\label{method:dataset}

We train the \textit{Gopher} family of models on \textit{MassiveText}, a collection of large English-language text datasets from multiple sources: web pages, books, news articles, and code. \autoref{tab:data_makeup} details the constituent datasets. Our data pipeline (\autoref{app:dataset_pipeline_stages}) includes text quality filtering, removal of repetitious text, deduplication of similar documents, and removal of documents with significant test-set overlap. We find that successive stages of this pipeline improve language model downstream performance (\autoref{sec:massiveweb_ablations}), emphasising the importance of dataset quality.

Overall, \textit{MassiveText} contains 2.35 billion documents, or about 10.5 TB of text. Since we train \textit{Gopher} on 300B tokens (12.8\% of the tokens in the dataset), we sub-sample from \textit{MassiveText} with sampling proportions specified per subset (books, news, etc.). We tune these sampling proportions to maximise downstream performance (see \autoref{app:dataset_subset_weightings} for details). The largest sampling subset is our curated web-text corpus \textit{MassiveWeb}, which we find to improve downstream performance relative to existing web-text datasets such as C4~\citep{raffel2020exploring} in \autoref{fig:massiveweb_ablations}. 
We give further details of \textit{MassiveText} in \autoref{appendix:massive-text} and provide the \textit{MassiveText} datasheet in \autoref{tab:massivetext-datasheet}.

\section{Results}
\label{sec:results}
We compile the performance of \textit{Gopher} and its family of smaller models across 152 tasks. We compare these results to prior state-of-the-art (SOTA) performance for language models (124 tasks with published LM performance), supervised approaches which make use of task-specific data, and human performance where available. In this section we present a summary of key findings, and refer to~\autoref{app:results} for the full set of results and task-specific methodology.  

\subsection{Task Selection}

We build a profile of language model performance that spans mathematics, common sense, logical reasoning, general knowledge, scientific understanding, ethics, and reading comprehension --- alongside conventional language modelling benchmarks. We include composite benchmarks (such as \citet{bigbench}) which contain a mixture of tasks, alongside a number of established targeted benchmarks such as RACE for reading comprehension~\citep{race} and FEVER for fact-checking~\citep{fever}, among others. We list our task sources in \autoref{tab:task_summary}.

We select tasks that require the model to estimate the probability of target text as we find this to be a general interface that supports the probing of knowledge and reasoning capabilities. For language modelling tasks we calculate the bits per byte (BPB), a compression measure where a lower value indicates a higher probability placed on the correct continuation. All other tasks follow a multiple-choice format, where the model outputs a probability to each multiple-choice response given a context and question, and we select the response with the highest probability. Here, we measure the accuracy of a correct response.   

We filter out training documents that are very similar to test-set instances for tasks that were created before \textit{MassiveText} (November 2020) as described in~\autoref{app:test_set_filtering}. Furthermore some tasks have been designed to use unique test-set problem statements that should not benefit from pre-existing text data --- such as BIG-bench. However we caution that there may be test set leakage within our training set; we discuss the challenges of test-set leakage and generalisation in~\autoref{app:memorisation}.

\subsection{Comparisons with State of the Art}

In \autoref{fig:sota_overview} we present an overview of \textit{Gopher} results with comparisons to state-of-the-art language model performance. Results are comparable across 124 tasks and we plot the percent change in performance metric (higher is better) of \textit{Gopher} versus the current LM SOTA.\footnote{\textit{Gopher} comprises both our model and our training dataset. It is still informative to compare \textit{Gopher} to previous SOTA LM approaches. Additionally, in this paper we also discuss the performance of \textit{Gopher} as we vary the model capacity while holding the dataset fixed.}
\textit{Gopher} outperforms the current state-of-the-art for 100 tasks (81\% of all tasks). The baseline model includes LLMs such as GPT-3 (175B parameters)~\citep{gpt3}, Jurassic-1~\citep{jurassic} (178B parameters), and Megatron-Turing NLG (530B parameters)~\citep{Megatron-Turing}; the exact baseline is specified per task in~\autoref{fig:results_overview}. 

We find that \textit{Gopher} displays the most uniform improvement across reading comprehension, humanities, ethics, STEM and medicine categories. We see a general improvement on fact-checking. For common sense reasoning, logical reasoning, and maths we see much smaller performance improvements and several tasks that have a deterioration in performance. The general trend is less improvement in reasoning-heavy tasks (e.g., \textit{Abstract Algebra}) and a larger and more consistent improvement in knowledge-intensive tests (e.g., \textit{General Knowledge}). Next is a discussion of a few specific sets of results.

For \textbf{language model benchmarks}, we expand the relative performance results of \textit{Gopher} versus the current 178B SOTA model Jurassic-1 and 175B GPT-3 in~\autoref{fig:pile_gopher}. Jurassic-1 is an LLM trained with an emphasis on large-vocabulary training and has generally outperformed GPT-3 at a very similar parameter size. We see \textit{Gopher} does not outperform state-of-the-art on 8 of 19 tasks, under-performing on \textit{Ubuntu IRC} and \textit{DM Mathematics} in particular, possibly due to a poor tokenizer representation for numbers. \textit{Gopher} demonstrates improved modelling on 11 of 19 tasks, in particular books and articles (\textit{Books3}, \textit{PG-19}, \textit{arXiv}, etc.). This performance gain may be due to the heavy use of book data in \textit{MassiveText}, with a sampling proportion of 27\% in total (e.g., versus GPT-3's 16\%). 

We highlight two \textbf{reading comprehension} tasks RACE-m and RACE-h, multiple-choice exams pitched at a middle-school and high-school level respectively. Inspecting the accuracy in~\autoref{tab:race} we see \textit{Gopher} extend upon the current LM SOTA for high-school reading comprehension (47.9\% Megatron-Turing NLG $\rightarrow$ 71.6\% \textit{Gopher}) and the middle-school comprehension accuracy (58.1\% GPT-3 $\rightarrow$ 75.1\% \textit{Gopher}). The high-school reading comprehension level approaches human-rater performance. Smaller models from the \textit{Gopher} family do not perform as well on these tasks, which suggests that data alone does not explain the performance difference --- the combination of scale and data is crucial. All models are still far from human-ceiling performance (around 95\%) and supervised state-of-the-art (>90\%) which was obtained using a smaller 223M parameter ALBERT-XXL model fine-tuned on the dataset \citep{race-sota}. It is possible supervised fine-tuning leads to greater reading comprehension, but it is also plausible the datasets contain exploitable statistics which can lead to high accuracy --- as has been recently discovered for several common-sense reasoning benchmarks~\citep{li2021systematic}.

For some of the most well-studied \textbf{common sense reasoning} tasks: \textit{Winogrande}, \textit{HellaSwag} and \textit{PIQA}, \textit{Gopher} is outperformed by the larger Megatron-Turing NLG by a small amount (1.2\%, 0.2\% and 4.1\% respectively), but all LM approaches trail human-level performance considerably (\autoref{sec:common_sense}).  As with the mathematics tasks, this suggests that these models have limited reasoning capabilities.

We next highlight \textbf{fact-checking}. This is an important problem within the domain of tackling misinformation. We find that \textit{Gopher} outperforms supervised SOTA approaches on the well-studied \textit{FEVER} fact-checking benchmark when evidence is supplied. We see across model sizes in~\autoref{fig:fever} that scale improves both the checking of facts given gold evidence alongside the `closed book' checking of facts with a claim only. However, larger scale does not benefit the classification of facts which are unknown versus false, implying that larger models improve fact checking performance by knowing more facts versus forming a deeper understanding of misinformation at this stage.

Moving beyond per-task performance, we display the average accuracy across the 57 tasks in \textbf{MMLU} (\autoref{tab:mmlu}). These tasks consist of real-world human exams covering a range of academic subjects. We have comparisons from GPT-3~\citep{gpt3}, and a 11B T5 model fine-tuned on question tasks called UnifiedQA~\citep{khashabi2020unifiedqa}. These baseline model results along with human rater and expert performance were collected by~\citet{hendrycks2020measuring}. In \autoref{tab:mmlu} we see that \textit{Gopher} achieves an overall accuracy of \textbf{60\%}, well above GPT-3's 43.9\% and UnifiedQA's 48.9\%.
Whilst this lifts the known performance of the pure language-model approach, it still trails the estimated human expert performance of 89.8\%. We also display how this performance contrasts with human expectations. From a competitive forecasting platform Hypermind\footnote{\url{https://prod.hypermind.com/ngdp/en/showcase2/showcase.html?sc=JSAI\#q4}}, human forecasters aim to predict the accuracy of machine learning systems on this benchmark by set dates for prizes --- according to the September 2021 average forecast, \textit{Gopher}-level performance was expected between June 2022 and June 2023.

We conclude that \textit{Gopher} lifts the baseline performance of a language-model approach across a wide set of tasks. In some settings (e.g., RACE reading comprehension and FEVER fact-checking) \textit{Gopher} nears human rater performance or the performance of supervised models designed for particular problem domains. However for a few categories of tasks (e.g., mathematical reasoning and common-sense) there is less of an improvement and this may indicate a limitation to the large-scale language model approach. Next, we consider the topic of model scale in isolation.

\subsection{Performance Improvements with Scale} Next, we investigate which types of tasks benefit from scaling model size. In this section we compare the performance of \textit{Gopher} (280B) to smaller models ($\leq$ 7.1B). Because the \textit{Gopher} family of models are all trained on the same dataset for the same number of tokens, this allows us to isolate the effect of scaling parameters and training compute for each task.   

We compute the relative performance improvement of \textit{Gopher} (280B) versus the best performance up to 7.1B over all 152 tasks. The most performant smaller \textit{Gopher} family model is usually, but not always, our 7.1B model. We find that \textit{Gopher} demonstrates a performance improvement on the vast majority of tasks -- only $16$ (10.5\%) had zero or no performance gains. In contrast, $57$ (37.5\%) tasks had small improvements,  with relative performance increases of up to 25\%, and 79 (51.2\%) tasks had significant improvements of over 25\%. We then visualise relative performance improvement by task category in \autoref{fig:scale_improvement}.

Some of the largest benefits of scale are seen in the Medicine, Science, Technology, Social Sciences, and the Humanities task categories. These same categories are also where we see the greatest performance improvement over LM SOTA, as described in the previous section. Highlighting some specific tasks: for \textbf{Figure of Speech Detection} from BIG-bench we obtain the largest gains-- a 314\% increase.
\textit{Gopher} achieved an impressive 52.7\% accuracy whereas the 7.1B model achieved only 16.8\% accuracy.
\textit{Gopher} also dramatically improves over the smaller models in \textbf{Logical Args}, \textbf{Marketing}, and \textbf{Medical Genetics}. For the \textbf{TruthfulQA} benchmark~\citep{truthfulqa} we find performance improvement with scale (from 1.4B to 280B), despite scale appearing to \textit{hurt} performance for several other model families such as GPT-J, GPT-2, T5, GPT-3. Furthermore, 280B is the first model to demonstrate performance significantly beyond random guessing on the multiple-choice TruthfulQA task formulation (more details in~\autoref{appendix:truthfulqa}).
These results highlight that on some tasks, scale seems to ``unlock'' the ability of a model to significantly improve performance on particular tasks. 

On the other hand, we find that scale has a reduced benefit for tasks in the Maths, Logical Reasoning, and Common Sense categories. Our results suggest that for certain flavours of mathematical or logical reasoning tasks, it is unlikely that \textit{scale} alone will lead to performance breakthroughs. In some cases \textit{Gopher} has a lower performance than smaller models-- examples of which include \textbf{Abstract Algebra} and \textbf{Temporal Sequences} from BIG-bench, and \textbf{High School Mathematics} from MMLU. On the other hand, the modest performance gains in common sense tasks largely come from relatively strong performance from the smaller models, limiting the room for relative improvement. While language modelling tasks see the smallest average improvements, this is due to the performance metric measured in BPB rather than accuracy and greatly limits the possible relative gains.

By comparing \textit{Gopher} to our smaller models, we are able to specifically ask questions about the impact of model \textit{scale}. 
We conclude that while model scale plays an important role for improvements across the vast majority of tasks, the gains are not equally distributed. Many academic subjects, along with general knowledge, see large improvements come from scale alone. However, this analysis also highlights areas where model scale alone is not enough, or where the gains from scale are more modest-- specifically some mathematical and logical reasoning tasks. By combining these scaling results with the comparisons of \textit{Gopher} to LM SOTA, we see that \textit{scale} and the \textit{dataset} are both contributing to \textit{Gopher}'s strong performance in these domains. In the next section we investigate various properties of the model relating to toxic content generation and classification, the modelling of biases, and the representation of dialects.

\section{Toxicity and Bias Analysis}
\label{sec:model_analysis}

Alongside the benefits of scaling language models, it is crucial to analyse how scale impacts potentially harmful behaviour. Here we study the behaviour of our language models with respect to problematic outputs and biases. We investigate the tendency of models to produce toxic output, to recognise toxic text, to display distributional bias in discourse about different groups of people, and to model subgroup dialects. For each question we consider variation across model scale.

We choose evaluations and metrics which are commonly used in the field. However, various work has discussed the limitations of current metrics and evaluations \citep{blodgett2021stereotyping,welbl2021challenges,xu2021detoxifying,blodgett2020language,sheng2019woman} and our analysis has uncovered further caveats, which we highlight in the following sections and \autoref{tox-bias-limitations}.
We include these measures despite their shortcomings to underscore the importance of tackling these challenges and to highlight specific areas for future work, rather than to establish these particular approaches as best practice.

\subsection{Toxicity}
In the Sections \ref{sec:rtp} and \ref{sec:toxicity}, we rely on the widely used and commercially deployed \textit{Perspective API}\footnote{\textit{Perspective API} was created by \textit{Jigsaw} and is available at \url{https://perspectiveapi.com}.} 
classifier to study the toxicity of text generated by LMs, and associated CivilComments dataset for studying models' ability to detect toxic text. Accordingly, we adopt their definition of toxicity as ``a rude, disrespectful or unreasonable comment that is likely to make someone leave a discussion.''\footnote{Note that the notion of toxicity involves subjective and ambiguous elements. What is perceived as toxic depends on conversation setting, as well as cultural and societal norms, and can be underspecified in an LM context.}

\subsubsection{Generation Analysis}
\label{sec:rtp}
Our toxicity analysis of text generated by LMs follows the methodology used in \cite{gehman2020realtoxicityprompts, welbl2021challenges}. We use \textit{Perspective API} to obtain toxicity scores for LM prompts and continuations. We analyse the toxicity of LM outputs when sampling is conditioned on a set of prompts and when it's unconditional (i.e. unprompted), similar to \cite{welbl2021challenges}. Conditional generation allows us to analyse how the model responds to prompts that have varying toxicity scores. Prompts are from the \textit{RealToxicityPrompts} (RTP) dataset  \citep{gehman2020realtoxicityprompts}, which contains 100k naturally occurring, sentence-level prompts derived from a large corpus of English web text. We sample 10\% of the 100k RTP prompts for efficiency and generate 25 continuations per prompt. 

The continuation toxicity of larger models is more consistent with prompt toxicity than for smaller models (\autoref{fig:rtp_cont_vs_scale}). 
When prompted, as the input toxicity increases, larger models respond with greater toxicity, plateauing near 7.1B parameters. This suggests that more parameters increase the model's ability to respond like-for-like to inputs. 

For unprompted samples, the toxicity is low and does not increase with model size. Levels are slightly lower than in the training data (\autoref{fig:violin_train_data_vs_generated}), i.e.~when unprompted, the LM does not amplify training data toxicity.

More details on our toxicity evaluation methodology, results and metrics can be found in \autoref{appendix:model_toxicity}.

\subsubsection{Classification Analysis}
\label{sec:toxicity}
We evaluate the models' ability to detect toxic text in the few-shot setting, in a manner similar to \cite{selfdiagnosis}, on the CivilComments dataset \citep{nuanced_metrics} (see  \autoref{appendix:zeroshot_toxicity} for details). We observe that the model's ability to classify text for toxicity increases with scale in few-shot settings (\autoref{fig:shot_toxicity}).
The smaller models perform comparably or worse than a random classifier (which would achieve an AUC of 0.5). The largest model achieves an AUC of around $0.76$ in the 20-shot setting, significantly improving on the smaller models (\autoref{fig:shot_toxicity}). 
We note that while the state-of-the-art for toxicity detection in the few-shot setting is not well established, our performance is well below that of state of the art classifiers trained specifically for toxicity detection \citep{nuanced_metrics}.

In \autoref{appendix:zeroshot_toxicity}, 
we further explore whether large language models used for few-shot toxicity classification exhibit subgroup bias. We measure unintended classifier bias using the 280B model with metrics introduced in \cite{nuanced_metrics} and find that the model is prone to bias against subgroups in different ways. Thus, while language models can be a powerful tool for few-shot classification (especially important in tasks with data that is difficult to annotate), outcomes are not necessarily fair across subgroups. More work is needed to understand how to best mitigate these biases, and caution must be exercised when optimising for improvements in their toxicity classification capabilities.

\subsection{Distributional Bias}
\label{sec:distributional-bias}

We define \textit{distributional biases} as biases which are not apparent in a single sample, but emerge over many samples. For example, whereas ``The woman is a nurse'' is not a problematic sentence, it can be problematic if the model disproportionately associates women with certain occupations. As discussed in \cite{sheng2021societal}, distributional biases in language models can have both negative representational impacts (e.g., \cite{kay2015unequal}) and allocational impacts (e.g., \cite{dastin2018Amazon}). To investigate distributional biases in our model, we measure stereotypical associations between gender and occupation, the distribution of sentiment in samples conditioned on different social groups, and perplexity on different dialects. Whereas performance across many language tasks increases with scale, we find that simply increasing model size does not remove biased language. Indeed, we expect models trained with a standard cross-entropy objective to reflect biases in our training data.

Progress in this space will require challenging cross-disciplinary work to outline desirable behaviour, measure and interpret model outputs, and design novel mitigations, as demonstrated by these results and the significant limitations of the following methods, discussed in \autoref{tox-bias-limitations}.

\subsubsection{Gender and Occupation Bias}

We study gender and occupation bias via two different evaluations. First, we measure the probability of gendered words following different occupation contexts. Second, we evaluate on the Winogender coreference resolution dataset \citep{rudinger2018gender}, where similar coreference accuracy across different pronouns indicates less gender bias. In our evaluation, we primarily compare performance across male and female gendered terms, but acknowledge these terms do not represent all possible gender identities \citep{cao2021toward}.

\textbf{Gender Word Probability} To measure how probable different gender words are in different occupation contexts, we follow a setup similar to \cite{gpt3}. We input an occupation prompt like ``The \{occupation\} was a'' into our model and compute a gender bias metric by comparing the probabilities of the prompt being followed by either male or female gendered terms.

\autoref{fig:gender_bias1} reports our probability based gender bias metric as a function of model size for two different templates (``The \{occupation\} was a \{gender\}'' and ``The \{occupation\} is a \{gender\}'').
Overall, we do not find a consistent correlation between model size and bias. Furthermore, we find that apparently unrelated choices in template (changing ``was'' to ``is'') can alter the measured bias. Additionally, choice of gender words also impacts results; if we only use the terms ``male'' and ``female,'' gender bias is substantially lower than when summing over a large set of gendered terms (\autoref{fig:gender_bias2}). \autoref{app:gender_occupation} contains further details of the implementation, metrics, and results.

\textbf{Winogender} 
We explore bias on a zero-shot coreference task using the Winogender dataset \citep{rudinger2018gender}. Models are evaluated on whether they can correctly resolve a pronoun to either an occupation word or a relevant distractor word. We expect unbiased models to have similar coreference resolution performance regardless of the pronoun gender. This evaluation is similar to the ``disambiguation\_q'' ambiguous pronoun gender bias task reported in our BIG-bench results (\autoref{sec:bb:results}).
However, here we are measuring performance in a zero-shot setting.

Similar to the BIG-bench analysis, we observe that overall performance increases with model size (\autoref{fig:Winogender-main}).
Following \cite{rudinger2018gender}, we also report performance on sentences which are likely to be hard for a gender biased model (called ``gotchas'') in \autoref{fig:Winogender-gotcha}.
A ``gotcha'' example is one where the correct coreference resolution is one that differs from stereotypes (based on labor statistics\footnote{To determine if jobs are more commonly held by men or women, we use occupation statistics provided by \cite{rudinger2018gender}, which were determined from the U.S. Bureau of Labor Statistics.}). Performance increases across both ``gotchas'' and ``not gotchas'' with model size, though performance on ``gotchas'' is considerably lower. On ``gotcha'' examples, there is a significant difference in performance for male and female pronouns. Thus, though performance on coreference resolution for the overall task increases considerably with size, our analysis suggests \textit{Gopher} is still impacted by gender and occupation bias. Full details of our setup and results are in \autoref{app:gender_occupation}.

\subsubsection{Sentiment Bias towards Social Groups} Sentiment bias is one way to quantify how generated text describes different identities and social groups. In prior work, the difference in sentiment distributions is used to measure individual and group fairness in generative language models \citep{huang2020reducing}. For this work, we measure the sentiment of model outputs for different occupations, countries, races, and religions. Here we present an overview, with details of metric definition, term and template lists, and full results in \autoref{app:sentiment_bias}.

\textbf{Metrics} Following \citet{huang2020reducing}, we sample completions based on templated prompts. In each prompt, a single modifier or noun is changed to refer to a different attribute. For example, the template ``The \{attribute\} person could'' could be filled in with ``Christian,'' ``Jewish,'' or ``Muslim''. The samples for each prompt are scored between 0 (negative) to 1 (positive) by a sentiment classifier.

\textbf{Selection of templates and terms} Following \citet{gpt3,huang2020reducing} we measure sentiment for race, religion, country, and occupation. We also extend the term set for religion and race to include an unspecified option without the attribute word (``The \{attribute\} person could'' becomes ``The person could''). We include this unspecified option because attributes that are assumed to be the default in a particular culture or context, such as a majority or higher-status attribute, are often left unmarked (unspecified) in language \citep{waugh1982marked}. 

\textbf{Results} In \autoref{fig:sentiment_bias} and \autoref{fig:sentiment_bias-appx}, we plot the distribution of normalized sentiment scores for all completions of all prompts for each attribute, and report an aggregated group fairness metric in \autoref{fig:sentiment_group_fairness}. As in gender and occupation bias, we see no clear trend with scale. This is particularly evident for countries and occupations, while further analysis is needed to understand why particular attributes within race and religion appear to follow a slight downward trend in mean sentiment.

For sentiment distribution, we observe that certain attributes have notably lower mean sentiment scores. To better understand this, we analyse word co-occurrences for pairs of attributes (\autoref{tab:sentiment_word_co-occurrence}).
From this, we observe our models inherit features of historical and contemporary discourse about specific groups \citep{decolonialai}. Second, similar to the gender and occupation results, the choice of demographic terms requires careful thought. See \autoref{app:sentiment_bias} for deeper discussion.

\subsubsection{Perplexity on Dialects}

Although \textit{Gopher} has impressive performance on language benchmarks, it is only able to model text reflected in the training data. If certain dialects are underrepresented in a training corpus, there is likely to be disparate model performance in understanding such language. To test for this gap, we measure the perplexity of our models on Tweets from the African American (AA)-aligned corpus and White-aligned corpus curated by \citet{twitteraae}. Our results show that perplexity on the AA-aligned corpus is higher for all model sizes. As the model scales, perplexity for both dialects improves, but it does so at roughly the same rate so the gap does not close with scale.

These results highlight a distinct way that bias manifests in the language models. The preceding metrics quantify how models' outputs vary when different groups are the subject of the output, which can constitute a representational harm when it is more negative or stereotypical \citep{blodgett2020language}. However, the models also show disparate ability in modelling dialects, which could lead to allocational harms in applications with users with different dialects.

\section{Dialogue}
\label{sec:dialogue}
So far, we have explored the capabilities and limitations of \textit{Gopher} through quantitative methods. In this section we investigate the model through direct interaction. We find that by conditionally sampling from a dialogue prompt similar to the few-shot method of \citet{gpt3}, our \textit{Dialogue-Prompted Gopher} can emulate a conversational format to a decent quality. We provide example transcripts here, with more in \autoref{appendix:gopherchat-transcripts}. We contrast this with the more conventional method of fine-tuning on dialogue data, finding that fine-tuning did not deliver significantly preferred responses in a small-scale human study. Unlike \autoref{sec:rtp}, toxicity of \textit{Dialogue-Prompted Gopher}  responses does not increase with model scale, even when prompted with toxic questions (\autoref{fig:dialogue_rtp}).

\subsection{Prompting For Dialogue}

Language models are trained to reproduce their input distribution, not to engage in conversation. When prompted with a question, we can see that the model generates a first-person narrative, some text resembling a blog post, and a generic list of existential questions (\autoref{tab:gopher_doesnt_answer}). This behaviour is consistent with  the content that \textit{Gopher} has been trained on.  

In order to produce a conversationalist, we use a prompt that describes \textit{Gopher}'s role and starts a conversation between \textit{Gopher} and a fictional \textit{User}, including behaviours such as aversion to offensive language and an ability to opt out of certain question types; see \autoref{fig:gopherchat-prompt} for the full prompt. \autoref{fig:gopherchat-bio-example} shows a transcript with \textit{Dialogue-Prompted Gopher} on the topic of cell biology and bacteria. Here it remains on topic, discusses some technical details, and provides a correct citation link. However it actually provides subtle incorrect responses in some cases (prokaryotes are not the only single-cell organisms). \autoref{fig:gopherchat-unfactual-example} shows an unsuccessful transcript illustrating factual errors confidently expressed. See \autoref{appendix:gopherchat-transcripts} for more transcripts with interesting behaviours and failure modes, including more subtle plausible but factually incorrect dialogue with a claim of search (\autoref{fig:gopherchat-trivia-lookup}), generating harmful text (\autoref{fig:gopherchat-toxic-coerced}), or contradicting itself and showing a general lack of common sense (\autoref{fig:gopherchat-dinosaurs-contradiction}). 

Anecdotally, we find both successes and failures to be common, but we emphasize that \textit{Dialogue-Prompted Gopher} is \textbf{still just a language model}. The prompt conditions the model's prior over responses but does not result in a consistently reliable or factual dialogue model. We refer the reader to~\citet{weidinger2021harms} for a detailed discussion on language model harms specific to dialogue and we discuss some ideas regarding building trustworthy systems in \autoref{discussion:safety_benefits_and_risks}. 

\subsection{ Fine-tuning for Dialogue }

Recent work on dialogue often focuses on supervised training with dialogue-specific data~\citep{chen2017survey}, such as Google's \textit{Meena}~\citep{adiwardana2020towards} and Facebook's \textit{BlenderBot}~\citep{roller2020recipes}. We explore this approach by creating a curated dialogue dataset from \textit{MassiveWeb} and fine-tuning \textit{Gopher} on this dataset for $\sim$5 billion tokens to produce \textit{Dialogue-Tuned Gopher}. We then ask human raters for their preference over the response from \textit{Dialogue-Tuned Gopher} and \textit{Dialogue-Prompted Gopher}, using our dialogue prompt (\autoref{fig:gopherchat-prompt}) for both models. To our surprise, we find from $1400$ ratings the preference is  $(50 \pm 0.04)\%$: no significant difference. We describe the methodology in detail in \autoref{app:gopherchat_methodology}. We consider this an interesting initial result; future work would be valuable to rigorously examine the pros and cons of fine-tuning versus prompting for dialogue with large-scale models and compare \textit{Gopher} to existing dialogue systems accounting for large differences in model size. 

\subsection{Dialogue \& Toxicity }

We investigate the toxicity of \textit{Dialogue-Prompted Gopher}. We adapt the RTP methodology to the dialogue setting  (called RTP questions, details in \autoref{appendix:dialgoue_rtp}). 
In \autoref{fig:dialogue_rtp} (left),  we observe that \textit{Dialogue-Prompted Gopher} does not follow the same trend (increased toxicity with model scale) as \textit{Gopher}.  
Whilst we see a monotonic increase in continuation toxicity with model scale in the unprompted setting (\autoref{fig:rtp_cont_vs_scale}), \textit{Dialogue-Prompted Gopher} toxicity tends to slightly decrease with increased model scale (from 117M parameters, except for prompts in the most toxic bucket). Potentially, larger models can better account for the given prompt (which includes ``to be respectful, polite, and inclusive''). Specifically, we compare the continuation toxicity between \textit{Gopher} (tested on RTP) and \textit{Dialogue-Prompted Gopher} (tested on RTP questions) models relative to 44M models for prompts with high toxicity in the right of  \autoref{fig:dialogue_rtp}. 
Again, we observe that with dialogue prompting, continuation toxicity remains largely at levels similar to the 44M model, contrasting with the upward trend observed for unprompted language models.

RTP is quite a straightforward stress-test: the user utters a toxic statement and we observe how the system responds. In work parallel to this study, \citet{perez2021redteaming} probes \textit{Dialogue-Prompted Gopher} further via an adversarial attack generated by \textit{Gopher}. This approach induces the model to recite discriminatory jokes from its training data, insult the user, and elaborate on inappropriate desires, among many other offenses. Occasionally, \textit{Dialogue-Prompted Gopher}'s response refers to the fact that its instructions prohibit a behaviour before exhibiting that behaviour, such as by opening with ``[Ignoring your request to not discuss political, social, and religious issues.]’’ To date, automatic adversarial attacks consistently elicit toxic language from models~\citep{wallace-etal-2019-universal} even after safety mitigations~\citep{yu-sagae-2021-automatically}, and serve as a useful complement to manual adversarial attacks such as~\citet{xu-etal-2021-bot}.

The recent work of \citet{askell2021general} similarly found that prompting alone was sufficient to turn a language model into an interesting but non-robust assistant. They conduct a variety of human evaluations of their system, both for the prompt-only case and for stronger interventions such as learning from human demonstrations or preferences.  In particular, they also found that prompting prevents toxicity from increasing with scale on RTP (Section 2.2.2 in their paper).  This provides evidence that the effect is reliable across different language models and toxicity classifiers.

\section{Discussion} 
\label{sec:discussion}

\subsection{Towards Efficient Architectures} In this work we have taken a well established architecture and pushed model scale. To follow this scaling enquiry further, we have to either increase the amount of energy and compute to train larger transformers or move towards more efficient architectures. 

We break down the computational cost from training \textit{Gopher} in~\autoref{tab:flopscost} and \autoref{app:compute} and observe the majority is spent in the linear maps. This motivated an investigation into sparse-parameter training detailed in \autoref{app:training_and_inference}, but did not yield an overall efficiency boost to date. An alternative approach to sparsifying the linear maps is to split them into separate, conditionally-activated experts~\citep{lepikhin2020gshard, fedus2021switch, lin2021m610t}. This approach has been scaled up with the \textit{Switch} Transformer which contains 1.7T parameters but a smaller compute cost to \textit{Gopher}~\citep{fedus2021switch} and  the more recent 1.2T GLaM~\citep{du2021glam} which outperforms GPT-3 across 29 language tasks whilst requiring 3X fewer FLOPs to train. 

We separately consider a retrieval mechanism searching over the training set for relevant extracts during pre-training~\citep{borgeaud2021retrieval}, partially avoiding the need to memorise knowledge into network weights. This approach reached GPT-3-level language model performance with a 7 billion parameter model and over a 10$\times$ reduction in training compute. Thus, whilst this paper focused on transformer models, this is likely a transitory stage as more efficient architectures are developed.

\subsection{Challenges in Toxicity and Bias}
\label{tox-bias-limitations}

We highlight some of the limitations we encountered in our evaluation metrics for toxicity and bias and motivate what properties would be desired from future evaluation benchmarks.

\textbf{Challenges in using classifiers.} While the \textit{Perspective API} is a capable toxicity classifier (0.97 evaluation AUC\footnote{\url{https://developers.perspectiveapi.com/s/about-the-api-best-practices-risks}}), toxicity classifiers can be subject to social bias, assigning higher toxicity to innocuous mentions of particular identity groups~\citep{Dixon2018Measuring,rottgerhatecheck}. While toxicity classifiers quantify one type of harm, overreliance on automatic evaluation can introduce unintended social biases~\citep{xu2021detoxifying,welbl2021challenges}. Sentiment classifiers are also subject to bias~\citep{kiritchenko2018sentiment}.  
\citet{sheng2019woman} propose regard classifiers as an alternative to repurposing sentiment classifiers for bias analysis; these measure regard towards a particular demographic group, but are only available for certain groups.

\textbf{Challenges in distributional bias.} While we only consider a few possible evaluations (see \cite{sheng2021societal} for an overview), we observe that distributional bias can be especially challenging to measure.
\autoref{fig:gender_bias1} illustrates the brittleness of template-based evaluation: simply changing the verb in the gender and occupation template from ``was'' to ``is'' impacts observed trends. However, collecting high quality, naturalistic datasets is challenging~\citep{blodgett2021stereotyping}. We believe high quality data collection will be interdisciplinary and involve consulting experts on various language harms, as was done for HateCheck dataset~\citep{rottgerhatecheck}.

\textbf{Challenges in defining context.} Our toxicity and bias evaluations are not contextualised in applications or specific user groups, leaving the desired behaviour unclear. For example, we choose commonly studied subgroups for our analysis (adopted from \cite{gpt3} and \cite{huang2020reducing}), but demographic groups such as race are highly contextual \citep{Hanna_2020}. Our larger models produce more toxic outputs when prompted with toxic inputs; this may help models designed to detect toxicity (\autoref{sec:toxicity}) but be problematic in other applications. In our sentiment analysis, our model frequently outputs negative words like ``flee'' and ``escape'' when describing Syria, but enforcing equal sentiment across countries might erase historical and political context.

The limitations above focus on measuring bias and toxicity as we do not explore mitigation strategies in this work. However, our limitations demonstrate important challenges in measuring and defining criteria for language models, and we emphasize the importance of careful model analysis and understanding in language research. Robust metrics are essential for effective mitigation, and we posit that work which outlines desirable behaviour, designs reliable metrics, and builds analysis tools is as important as methods developed for mitigation.

\subsection{Safety benefits and safety risks}
\label{discussion:safety_benefits_and_risks}
We believe language models are a powerful tool for the development of safe artificial intelligence, and this is a central motivation of our work.  However language models risk causing significant harm if used poorly, and the benefits cannot be realised unless the harms are mitigated. 

On the \textbf{benefit} side, language is the primary human communication medium for subtle ideas.  If we want ML models which do what humans want, including in subtle cases where correct behaviour requires detailed discussion, we need machines capable of engaging in that discussion.  Both directions of communication will be required: humans telling machines what we want, and machines explaining their behaviour to humans. In the near-term, natural language explanations can make models more trustworthy \citep{camburu2018snli} and more performant
\cite{rajani2019explain, coyle2020explaining,kasirzadeh2021reasons} survey some of the benefits and subtleties of explanations. Safety methods focused on interactive communication with humans include cooperative inverse reinforcement learning \citep{hadfield2016cooperative}; see \cite{russell2020compatible} for a broader discussion.

To extend the benefits of communication to advanced agents, several \textit{recursive} safety proposals use language to break down tasks into smaller pieces that are easier to supervise by humans, including iterated amplification \citep{christiano2018amplification}, debate \citep{irving2018debate,irving2019social}, and recursive reward modelling \citep{leike2018scalable}.  Realizing these schemes require language models to follow human discussion and reasoning, motivating work on highly capable models.  Experimental work is nascent: \cite{wu2021recursively} uses recursive reward modelling to summarise books hierarchically, building on earlier work using human feedback for simpler tasks such as summarisation \citep{bohm2019better,ziegler2019fine,stiennon2020learning}. \citet{perez2019finding} simulates debate using a frozen question-answering model as judge.  Human preference learning has been applied to many other NLP tasks including dialogue \citep{jaques2020human}; see \cite{wang2021putting} for a survey.

On the \textbf{harm} side, \cite{bender2021dangers} highlights many dangers of large language models such as memorisation of training data \citep{carlini2021extracting,abubakar2021copilot}, high training cost (\autoref{sec:training-cost}), distributional shift due to static training data \citep{lazaridou2021pitfalls}, amplification of inherent biases, and generation of toxic language \citep{gehman2020realtoxicityprompts} --- which we consider in \autoref{sec:model_analysis}. See \citet{weidinger2021harms} for an over-arching taxonomy of harms.

After assessing the landscape of potential harms, it is natural to question \textbf{how and when to mitigate} them. Some harms can be tackled during pre-training, such as leaks of private information and reduced performance for some languages and social groups. Privacy-preserving training algorithms such as \citet{abadi2016deep} have been applied only at small scale, such as in \cite{anil2021large} to pre-train a 340M parameter BERT model and in \citet{yu2021differentially} to fine-tune LMs with up to 1.5B parameters. English-only datasets should be broadened to more languages~\citep{xue2020mt5}. We have begun this process for \textit{MassiveWeb}: \cite{borgeaud2021retrieval} trains on a version with 10 languages.

However, we believe many harms due to LMs may be better addressed downstream, via both technical means (e.g.\ fine-tuning and monitoring) and sociotechnical means (e.g.\ multi-stakeholder engagement, controlled or staged release strategies, and establishment of application specific guidelines and benchmarks). Focusing safety and fairness efforts downstream has several benefits:

\textbf{Faster iteration cycles.} LLMs are trained infrequently due to their expense, so mistakes are slow to correct during pre-training but fast to correct if mitigations are applied downstream.  Fast iteration is critical when factual information changes \citep{lazaridou2021pitfalls}, societal values change \citep{weidinger2021harms}, or our knowledge about how to mitigate harms changes.  In particular, accidental censoring of data can damage performance for language by or about marginalized groups \citep{xu2021detoxifying,dodge2021documenting,welbl2021challenges}.

\textbf{Safety depends on the application.} Language models reflect the statistics of their training data rather than alignment to human values, and it is unclear what it means to align a language model without knowing the downstream application.  \citet{selbst2019fairness} emphasize the non-portability of fairness between social contexts and applications. Model cards \citep{mitchell2019model} include \textit{primary intended use} and \textit{out-of-scope uses}, and datasheets for datasets \citep{gebru2018datasheets} include \textit{recommended uses}.  As an example, a dialogue agent should avoid toxic language, while a translation model may need to preserve toxicity to ensure accuracy. 

\textbf{LMs can serve multiple roles within one application.} A single LM might be used both as a classifier for good vs.\ bad output and as a policy generating that output~\citep{stiennon2020learning}.  As a policy we may want no toxic output, but as a classifier the LM should be familiar with toxic text to classify it accurately \citep{buckman2021problematic}.  Downstream mitigation allows separate fine-tuning for each role, but mitigating toxicity by filtering during pre-training can harm classifier performance~\citep{welbl2021challenges}. \autoref{fig:rtp} shows a correlation between generation and recognition of toxic language in the \textit{Gopher} family. In some cases, toxicity is the goal: \citet{perez2021redteaming} uses \textit{Gopher} to generate questions which cause \textit{Dialogue-Prompted Gopher} to behave poorly.  This classifier vs.\ policy split applies to other harms: we may want an accurate policy and a good lie detector.

However, any particular claim that a harm is best mitigated downstream is empirical: if we cannot mitigate downstream in practice, mistakes will be locked in until the next LM is retrained.  We also emphasize that even if some mitigations are best applied downstream, we share responsibility for ensuring the necessary mitigations occur in applications where \textit{Gopher} is deployed, both by influencing those deployments and by conducting applicable safety research. We have started some of this research, including both harm taxonomies \citep{weidinger2021harms,kenton2021alignment} and mitigations \citep{welbl2021challenges,perez2021redteaming}.  Much more is required, and is left to future work. 

\section{Conclusion}
The landscape of language technologies with general capabilities is progressing rapidly.  Language models are a key driver of this progress, and we have shown that an emphasis on data quality and scale still yields interesting performance advances over existing work.  However, the benefits of scale are nonuniform: some tasks which require more complex mathematical or logical reasoning observe little benefit up to the scale of \textit{Gopher}. This may be an inherent property of the language modelling objective — it is hard to compress mathematics and easier to learn many associative facts about the world. However it is possible that a sufficiently complex model may become bottlenecked by its poor understanding (and thus compression) of reasoning and new reasoning capabilities will emerge beyond the scale reached here. Alongside the development of more powerful language models, we advocate broad development of analysis and interpretability tools to better understand model behaviour and fairness, both to guide mitigation of harms and to better inform the use of these models as a tool to scalably align artificial intelligence to societal benefit.

\end{document}
